\textbf{Specialized propagation} is a propagational algorithm based on the \textbf{semantic} of the constraint. So, we don't put its propagation in a naive way but 
we do it in the fastest way exploiting the semantic of that specific constraint. \vspace{7pt}

Such approach can be potentially much more efficient than a general approach, but it's not always easy to develop. \vspace{7pt} 

As we already know, \textbf{global constraints} are very useful for the modelling phase, but at the same time they embed specialized propagation algorithms
making the propagational process much more effective than the other ways to express the same complex constraints. \vspace{7pt}

We can see global constraints benefits in two perspectives:
\begin{itemize}
    \renewcommand{\labelitemi}{-}
    \item \textbf{Modelling benefits}. 
    \begin{itemize}
        \item Reduce the gap between the problem description and the formal description (\textit{the model}).
        \item May allow the expression of constraints that are otherwise not possible to state using primitive constraints.
    \end{itemize} 
    \item \textbf{Solving benefits}. 
    \begin{itemize}
        \item Efficient propagation.
        \item Strong inference in propagation.
    \end{itemize} 
\end{itemize}
Therefore, as rule of thumb, any time we have some complex relationships to express, we should check the existing global constraints and see if we can exploit 
any of them.  